\section{稳定性与性能分析}

本节在理想连续时间条件下给出闭环预设时间跟踪结论，并进一步说明模型误差、补偿误差和数值实现对该结论的影响。除特别说明外，假设观测器在 \(T_o\) 内完成收敛，且 PDF 控制器从 \(t=T_o\) 起以局部时间 \(\vartheta=t-T_o\) 启动。

\textbf{命题1：} 若预设时间扰动观测器满足其连续时间收敛条件，且 \(\Delta_a\) 与 \(\dot\Delta_a\) 有界，则存在设计时间 \(T_o>0\)，使得
\[
\hat\Delta_a(t)=\Delta_a(t),\qquad \hat d(t)=d(t),
\qquad \forall t\ge T_o.
\]

\textbf{说明：} 命题1是连续时间理想观测器结论。由
\[
\hat d=M_e(q_m)\hat\Delta_a,
\qquad
\Delta_a=M_e^{-1}(q_m)d
\]
可知，当 \(\hat\Delta_a=\Delta_a\) 时有 \(\hat d=d\)。在采样实现、测量噪声或观测器公式未完全满足理论条件时，不能直接将该结论解释为数值观测误差严格等于零。

\textbf{命题2：} 假设命题1成立。若 \((A,B)\) 可控，\(K_0\) 使
\[
A_c=A+BK_0
\]
为期望的基准闭环矩阵，并且周期增益 \(K_c(t)\) 按式 \eqref{eq:kc_pdf} 构造，使相应单周期状态转移矩阵满足
\[
\Delta_K(h)=0,
\]
则理想模型条件下，关节跟踪误差满足
\[
e(t)=0,\qquad \dot e(t)=0,
\qquad \forall t\ge T_o+2h.
\]

\textbf{证明：} 当 \(t\ge T_o\) 时，由命题1有 \(\hat d=d\)。将控制律 \eqref{eq:computed_torque} 代入式 \eqref{eq:joint_model}，得到 \(\ddot q_m=v\)。再令
\[
v=\ddot q_d+\nu,
\]
则有 \(\ddot e=\nu\)，并且
\[
\dot z=Az+B\nu,
\qquad
A=\begin{bmatrix}0&I_n\\0&0\end{bmatrix},
\quad
B=\begin{bmatrix}0\\I_n\end{bmatrix}.
\]
代入 \(\nu=K_0z-K_c(t)z(t-h)\)，得
\[
\dot z=A_cz-BK_c(t)z(t-h).
\]
在 PDF 周期的前半段 \(K_c(t)=0\)，因此可按分步法求得一个完整周期的状态映射。依据式 \eqref{eq:period_map} 以及式 \eqref{eq:kc_pdf} 的 Gramian 构造，闭环单周期状态转移矩阵满足
\[
\Delta_K(h)=0.
\]
因此局部时间达到 \(2h\) 后有 \(z(t)=0\)。由于局部时间为 \(\vartheta=t-T_o\)，故 \(t\ge T_o+2h\) 时 \(e(t)=0\) 且 \(\dot e(t)=0\)。\hfill\(\square\)

上述命题给出本文控制器的理想预设时间上界
\begin{equation}
T_{\mathrm{pre}}=T_o+2h.
\label{eq:prescribed_total_time}
\end{equation}
这里的“预设”体现在 \(T_o\) 和 \(h\) 可在控制器设计阶段选定，而不是由初始误差大小决定。需要注意，式 \eqref{eq:prescribed_total_time} 是连续时间理想结论；采样、力矩饱和和模型误差会使实际误差不再严格于该时刻归零。

\textbf{命题3：} 若采用一般周期延迟增益，并且相应单周期状态转移矩阵 \(\Delta(h)\) 为幂零矩阵，即存在 \(\nu_0\ge1\) 使
\[
\Delta^{\nu_0}(h)=0,
\]
则理想补偿条件下有
\[
e(t)=0,\qquad \dot e(t)=0,
\qquad \forall t\ge T_o+2\nu_0h.
\]
该结论来自周期状态映射
\[
z(2(k+1)h)=\Delta(h)z(2kh).
\]
式 \eqref{eq:kc_pdf} 对应更强的单周期零化情形，即 \(\nu_0=1\)。

\subsection{残余扰动下的实际跟踪性能}

当模型存在近似误差、扰动估计不完全或执行器受到饱和限制时，定义等效残余加速度扰动
\begin{equation}
\Delta_r(t)=M_e^{-1}(q_m)\bigl[d(t)-\hat d(t)+\Delta_m(t)+\Delta_\tau(t)\bigr],
\label{eq:residual_acceleration}
\end{equation}
其中 \(\Delta_m\) 表示名义模型与真实模型之间的等效动力学误差，\(\Delta_\tau\) 表示力矩指令经饱和或量化后产生的输入误差。此时误差系统可写为
\begin{equation}
\dot z(t)=A_cz(t)-BK_c(t)z(t-h)+B\Delta_r(t).
\label{eq:perturbed_delay_error}
\end{equation}
若 \(\Delta_r(t)\) 有界，则理想单周期零化映射被有界输入项取代，严格精确归零一般不再成立。下面给出一个基于单周期输入--状态映射的实际有界性结论。

定义 \(\mathcal S_h(\theta,s)\) 为齐次延迟系统
\[
\dot z(t)=A_cz(t)-BK_c(t)z(t-h)
\]
在一个 PDF 周期内从时刻 \(s\) 到时刻 \(\theta\) 的因果解算子，并定义周期内的扰动放大系数
\begin{equation}
\Gamma_h=\max_{0\le \theta\le 2h}
\int_0^\theta
\left\|\mathcal S_h(\theta,s)B\right\|\,\mathrm ds .
\label{eq:pdf_disturbance_gain}
\end{equation}
这里的范数可取诱导二范数。对于给定的 \(A_c\)、\(K_c\) 和 \(h\)，\(\Gamma_h\) 可通过离线数值积分或对相应延迟系统进行离散化计算。

\textbf{命题4：} 假设 PDF 增益满足理想单周期零化条件 \(\Delta_K(h)=0\)，且残余扰动满足
\[
\|\Delta_r\|_\infty
=\sup_{t\ge T_o}\|\Delta_r(t)\|\le\bar\Delta_r.
\]
则在连续时间实现下，完整 PDF 周期边界处的闭环误差满足
\begin{equation}
\|z(T_o+2(k+1)h)\|\le \Gamma_h\bar\Delta_r,
\qquad k=0,1,2,\ldots .
\label{eq:pdf_practical_bound}
\end{equation}
特别地，当 \(\Delta_r(t)\equiv0\) 时，式 \eqref{eq:pdf_practical_bound} 退化为理想的单周期精确归零结论。

\textbf{证明：} 令 \(z_k=z(T_o+2kh)\)，并将第 \(k\) 个 PDF 周期内的局部时间记为 \(\vartheta\in[0,2h]\)。变分公式给出
\[
z_{k+1}=\Delta_K(h)z_k+
\int_0^{2h}\mathcal S_h(2h,s)B\Delta_r(T_o+2kh+s)\,\mathrm ds.
\]
由 \(\Delta_K(h)=0\) 以及 \(\|\Delta_r\|_\infty\le\bar\Delta_r\)，有
\[
\|z_{k+1}\|
\le \int_0^{2h}
\|\mathcal S_h(2h,s)B\|\,\|\Delta_r(T_o+2kh+s)\|\,\mathrm ds
\le \Gamma_h\bar\Delta_r.
\]
该不等式对每个 \(k\ge0\) 成立，故得到式 \eqref{eq:pdf_practical_bound}。注意，该结论首先约束的是完整 PDF 周期的边界状态；若需要给出周期内任意时刻的误差界，还需额外定义周期内状态传播增益。\hfill\(\square\)

命题4给出了方向一所需的定量结论：完整 PDF 周期边界处的实际误差邻域大小由残余扰动上界 \(\bar\Delta_r\) 和 PDF 周期扰动增益 \(\Gamma_h\) 共同决定。若需要将结果推广到周期内任意时刻，还需补充周期内状态传播增益。进一步地，若数值实现或正则化计算只能保证
\[
\|\Delta_K(h)\|\le\varepsilon_\Delta<1,
\]
则在 PDF 周期边界处有递推不等式
\begin{equation}
\|z_{k+1}\|
\le \varepsilon_\Delta\|z_k\|
+\Gamma_h\bar\Delta_r,
\label{eq:pdf_approximate_recursion}
\end{equation}
从而
\begin{equation}
\limsup_{k\to\infty}\|z_k\|
\le \frac{\Gamma_h\bar\Delta_r}{1-\varepsilon_\Delta}.
\label{eq:pdf_ultimate_bound}
\end{equation}
式 \eqref{eq:pdf_ultimate_bound} 表明，Gramian 数值误差、采样误差和正则化误差会通过 \(\varepsilon_\Delta\) 放大残余扰动引起的稳态误差。由 \eqref{eq:residual_acceleration}，若
\[
\|M_e^{-1}(q_m)\|\le\bar m^{-1},\quad
\|d-\hat d\|\le\bar d_o,\quad
\|\Delta_m\|\le\bar\Delta_m,\quad
\|\Delta_\tau\|\le\bar\Delta_\tau,
\]
则可以采用保守估计
\begin{equation}
\bar\Delta_r\le\bar m^{-1}
\left(\bar d_o+\bar\Delta_m+\bar\Delta_\tau\right).
\label{eq:residual_bound_explicit}
\end{equation}
因此，理论误差界可进一步写成
\[
\|z(t)\|
\le \Gamma_h\bar m^{-1}
\left(\bar d_o+\bar\Delta_m+\bar\Delta_\tau\right)
\]
的形式；对于近似 PDF，则将右端替换为式 \eqref{eq:pdf_ultimate_bound} 给出的上界。

\subsection{抗饱和补偿器的稳定性分析}

当采用抗饱和补偿器 \eqref{eq:as_dynamics} 和修正控制律 \eqref{eq:final_controller_as} 时，式 \eqref{eq:residual_acceleration} 中的力矩输入误差 \(\Delta_\tau\) 由辅助变量 \(\eta\) 的动态补偿。此时残余扰动可重新写为
\begin{equation}
\Delta_r^{\mathrm{as}}(t)=M_e^{-1}(q_m)\bigl[d(t)-\hat d(t)+\Delta_m(t)\bigr]
+ \Delta_u(t) + \eta(t),
\label{eq:residual_as}
\end{equation}
其中 \(\Delta_u\) 与 \(\eta\) 通过式 \eqref{eq:as_dynamics} 关联。将 \eqref{eq:residual_as} 代入误差动态 \eqref{eq:perturbed_delay_error}，可得
\[
\dot z = A_c z - BK_c(t)z(t-h) + B(\eta + \Delta_u) + B\bar\Delta_{\mathrm{other}},
\]
其中 \(\bar\Delta_{\mathrm{other}}=M_e^{-1}(d-\hat d+\Delta_m)\) 为除饱和外的残余项。由 \(\eta\) 的动力学 \eqref{eq:as_dynamics} 可知，当系统处于饱和 \((\Delta_u\neq0)\) 时，\(\eta\) 累积饱和误差，并通过控制律 \eqref{eq:final_controller_as} 中的 \(\eta\) 项改变计算力矩，减小饱和深度；当饱和解除 \((\Delta_u=0)\) 后，\(\eta\) 在固定时间 \(T_{\eta,\max}\) 内收敛至零。因此，在所提出的抗饱和补偿器作用下，饱和引起的跟踪性能退化被限制在 \(\Delta_u\neq0\) 的时间段内，且饱和结束后系统可快速恢复至标称的 PDF 收敛特性。

因此，本文对理论结果作如下区分：
\begin{enumerate}
  \item 连续时间、精确模型、无饱和且观测器精确收敛时，得到式 \eqref{eq:prescribed_total_time} 所示的预设时间精确跟踪；
  \item 存在模型误差、持续扰动补偿误差或执行器非理想特性时，得到预设时间实际有界跟踪；
  \item 离散数值仿真结果用于验证控制结构和有限精度下的跟踪性能，不能单独作为连续时间严格精确归零定理的证明。
\end{enumerate}

对于末端任务空间轨迹，可通过广义雅可比矩阵生成关节参考速度
\[
\dot q_d=J_g^\#(q_m)\left[\dot x_e^d-K_x(x_e-x_e^d)\right],
\qquad
J_g^\#=J_g^\top(J_gJ_g^\top+\lambda^2I)^{-1}.
\]
该步骤属于上层参考生成器，而非本文 PDF 误差闭环的必要条件。若 \(J_g\) 接近奇异，还需结合阻尼调节、奇异规避、任务优先级和关节限位处理，以保证生成的 \(q_d,\dot q_d,\ddot q_d\) 满足问题描述中的有界性要求。

综上，严格结论依赖四个条件：扰动观测器在设计时间内精确收敛、动力学补偿足够准确、PDF 增益满足单周期零化或幂零条件，以及延迟历史能够稳定获得。任一条件被破坏时，应采用“预设时间实际有界”或“小邻域跟踪”的表述，而不应继续宣称严格精确归零。
